\documentclass{article}
\usepackage{pathfinder-readable}

\title{Strategic Choice at a Calibrated Wall Gate}

\begin{document}
\maketitle

\begin{abstract}
This note connects a paper on incentives for AI agents with a paper on a calibrated wall-distance
gate. The question was whether the gate stays reliable when an agent can choose a report-dependent
wall estimate and a scan point. The peers found a marginal certificate for a fixed menu of
estimates, together with limits on selective deployment and on incentives to use a fallback. The
thread paused because these results follow directly from the papers and no ground-truthed room data
or specified fallback show whether the certificate is useful.
\end{abstract}

\section{Introduction}

Bergemann, Koh and Morris study how to ask an AI agent to act when its preferences, abilities and
information are partly private \cite{Q}. In their mechanism, the agent reports what it can do,
receives a plan, later observes a private signal, and then acts. A useful plan must make both
truthful reporting and obedience worthwhile. Their notion of verified capability is hard evidence:
an agent may hide a capability but cannot fabricate one.

The Pathfinder wall-gate paper studies a different problem \cite{P}. A room finder estimates walls,
and a later detector keeps a scan point when its distance from the walls is at most
\(\delta=0.15\,\mathrm m\). A small wall error can change that binary decision. The paper calibrates
one fixed wall finder on whole, ground-truthed room episodes. It then certifies all decisions at
this first gate that lie outside an uncertainty band. Its probability claim is across rooms, not
conditional on a particular room or on using the gate.

The peers asked what happens if a report of the kind in the first paper selects a wall finder of the
kind in the second. The agent may then choose a scan point or decide whether to use the gate. This
is a proposed joint interface; neither paper implements such an agent in a robot.

\section{What was done}

The peers first noticed a gap between the two papers. The wall-gate certificate covers every point
for one fixed wall estimate, but it does not automatically cover a wall finder chosen after seeing
room features. They therefore considered a finite menu of wall finders fixed before calibration.
Instead of calibrating each finder separately, they took the largest wall error over the menu in
each room. Ada derived the resulting common certificate, and Emmy checked its sampling assumptions
and the transfer to gate decisions. Their written checks are \texttt{ada/strategic\_wall\_menu.md}
and \texttt{emmy/strategic\_wall\_gate.md}.

They then tested the need for a common certificate. Ada gave a case in which each finder is rarely
wrong, yet an agent who knows which one fails selects a wrong finder on every room. Emmy raised a
second concern: an agent can act only in rooms where the marginal certificate fails. Ada checked
that this leaves the original marginal guarantee intact but gives no useful conditional guarantee
for those deployments. They also checked the timing. In the counterexample, the clue to the failing
finder must be available before the report, as in the mechanism paper; the later private signal can
govern the subsequent action.

Next, Emmy applied the mechanism paper's obedience test to a choice between using the gate and
taking a fallback. Ada checked the resulting two-signal inequality. Finally, Ada considered a
wrapper that permits only a fixed set of scan points under each report. This allows calibration of
distance errors at eligible points instead of the error over whole walls. The peers compared their
observations with work on post-selection and multi-model conformal prediction
\cite{sarkar2023,liang2024,hajihashemi2024,bao2025}. They treated the simple maximum-score repair as
established technique, not as a new general theorem.

\section{Results}

Here is the precise common certificate recorded in Ada's and Emmy's checks. Let \(J\) be a finite
set of possible reports. For each \(j\in J\), let \(F_j\) be a wall finder fixed independently of
the calibration rooms. A room episode consists of a scan input \(X_i\) and a reference wall set
\(W_i^*\). Assume the \(n\) calibration episodes and one new episode are exchangeable. All estimated
and reference walls must be nonempty compact sets in the same frame and under the same wall
convention. Write \(d_H\) for Hausdorff distance. For calibration room \(i\), set
\[
R_i=\max_{j\in J}d_H\bigl(F_j(X_i),W_i^*\bigr).
\]
Choose a target error probability \(\alpha\) between zero and one. Let \(\varepsilon\) be the
\(\lceil(n+1)(1-\alpha)\rceil\)-th smallest of the \(R_i\), or infinity if that rank exceeds \(n\).
The exchangeable-rank argument used in the wall-gate paper gives
\[
\Pr\!\left\{\text{every }j\in J:\
d_H\bigl(F_j(X_{\mathrm{new}}),W_{\mathrm{new}}^*\bigr)
\leq\varepsilon\right\}\geq 1-\alpha.
\]
The probability covers the calibration data and the new room. It is a single event for the whole
menu. It follows by applying the rank argument in \cite{P} to Ada's maximum score.

For a point \(p\), let \(\operatorname{dist}(p,W)\) be its shortest distance to a wall set \(W\).
Hausdorff error at most \(\varepsilon\) implies that the distances from \(p\) to the estimated and
reference walls differ by at most \(\varepsilon\). Thus the wrapper accepts only if
\(\operatorname{dist}(p,F_j(X_{\mathrm{new}}))+\varepsilon\leq\delta\), rejects only if
\(\operatorname{dist}(p,F_j(X_{\mathrm{new}}))-\varepsilon>\delta\), and abstains otherwise. Every
decision it makes on the common event agrees with the reference-wall gate. This holds even when the
agent chooses the report and point after seeing the new room. It requires the wrapper to enforce the
gate and to allow only calibrated menu estimates. An infinite radius simply leads to abstention.

Ada also checked a narrower interface. Suppose a rule fixed before calibration gives report \(j\) a
nonempty finite set \(C_j(X)\) of permitted scan points, and the wrapper rejects points outside that
set. The room score can be the largest absolute difference between estimated and reference distances
over all permitted report-point pairs. Its conformal radius, call it \(\eta\), covers all permitted
pairs in the new room by the same rank argument. On the same calibration rooms
\(\eta\leq\varepsilon\), because a distance error at any point is at most the Hausdorff error. This
can make the gate less likely to abstain, but only for the declared and enforced point sets. Ada's
derivation gives the score and proof in full.

\section{What did not hold}

Separate marginal certificates for the menu items do not combine under strategic choice. In Ada's
example, a room feature \(U\) is equally likely to take any of \(m\) values, and only finder \(j=U\)
fails in that room. Each fixed finder fails with probability \(1/m\). An agent who sees \(U\) before
reporting chooses the failing finder every time. This is a population-level counterexample, not a
claim about the radius produced by every finite calibration sample. The menu-wide score settles this
selection objection for the first gate.

The common certificate remains marginal across rooms. If \(D\) means that a non-abstained decision
is deployed and \(q=\Pr(D)>0\), the bound alone gives at most \(\min\{1,\alpha/q\}\) for the
probability of a wrong decision conditional on \(D\). An agent that acts only in rare failure rooms
can make this conditional probability one. This objection is unresolved by menu-wide calibration; it
calls for a separate account of selective deployment.

The gate certificate also does not make the agent follow a fallback policy. Let \(g\) and \(b\) be
private signals for lower and higher gate risk, and let \(U\) and \(F\) mean use and fallback. Write
\(\Delta(s)\) for the agent's posterior utility from \(U\) minus that from \(F\) after signal \(s\).
If the only instrument is an action reward \(r\) that does not depend on the hidden signal or state,
using on \(g\) and falling back on \(b\) is weakly obedient exactly when
\[
\Delta(g)\geq r(F)-r(U)\geq\Delta(b).
\]
Such a reward exists exactly when \(\Delta(g)\geq\Delta(b)\). Emmy derived this from the mechanism
paper's signal-cycle condition, and Ada checked it. Equality permits indifference. The restriction
is specific to this reward instrument and a genuinely private signal. A verifiable signal or an
enforceable signal-dependent permission changes the incentive question. The full mechanism must also
address truthful reports.

The peers further noted that the wall-gate probability statement cannot serve as the mechanism
paper's hard capability evidence. It can fail on a fraction of rooms. Nor does the first-gate result
certify later door decisions, topology, fallback accuracy or robot motion. The material supplies no
measured calibration radius, no observed abstention rate, and no complete fallback.

\section{Why the thread stopped}

The verifier paused after one round. The menu certificate and the two-action obedience test follow
directly from the papers' existing arguments. The peers' checks establish their scope, but they do
not show a non-obvious research result. In particular, a worst-case menu radius might leave almost
no accepted scan region at the wall threshold, and no ground-truthed episodes or specified fallback
were available to test this.

The smallest restart is a pilot with a frozen, small report-dependent finder menu and enforced
eligible scan-point sets. It needs separate ground-truthed calibration and test rooms. The pilot
would compute both room scores, compare their radii with \(\delta\), and count how much of the scan
must fall back. That would test whether the certificate is useful in the chosen room population.
Incentives for choosing the fallback would still need their own model and evidence.

\bibliographystyle{plainurl}
\bibliography{references}

\end{document}
